\documentclass[11pt]{article}
\usepackage[a4paper,margin=1in]{geometry}
\usepackage{amsmath,amssymb}
\usepackage{graphicx}
\usepackage{booktabs}
\usepackage[hidelinks]{hyperref}
\usepackage{authblk}

\newcommand{\Fl}{\mathbb{F}_\ell}
\title{\textbf{Addendum: Two Adversarial Tests of the $\chi(3)$ Channel}\\
\large Frobenius traces and a level-6 automorphic form against the 6N wing}
\author{R.\ Chen}
\affil{GUT Geoservice Inc., Montr\'eal}
\date{\textit{6N Arithmetic Geodynamics series} $\cdot$ Addendum to the $\chi(3)$ Channel note $\cdot$ 2026}

\begin{document}
\maketitle

\noindent
The $\chi(3)$-channel principle asserts that the 6N wing of a prime $p$ (its class $p\bmod 3$) can
enter an $L$-function-derived invariant only through the $q=3$ Euler factor $\chi_p(3)$. A principle
of this kind earns its standing by surviving attempts to break it. This addendum reports two such
attempts---deliberately constructed to give the wing its best chance---together with the data that
defeat them. They fail by \emph{opposite} mechanisms and reach the \emph{same} verdict.

\section*{The two break-in attempts}
\textbf{Test A --- decrepitation (Galois fluid).} Bombard the curve $E_p : y^2 = x^3 - p$ at a prime
``temperature'' $\ell$ and extract the Frobenius trace $a_\ell(E_p)$, for $\ell = 7, 13$, over all
$2260$ primes $3 < p < 2\times10^4$. The hope: a quantity drawn from deep in the Galois representation
might carry wing information past the CRT lattice.

\noindent\textbf{Test B --- a native level-6 plume (Langlands).} Instead of a generic object, choose
one whose \emph{level} already carries $6$: the cusp-form newforms of $\Gamma_0(6)$, weights $k=4,6$
(the space $S_2(\Gamma_0(6))$ is $0$-dimensional, so $k=2$ is empty). Extract the Hecke eigenvalue
$a_p$ at all $667$ primes $3 < p < 5000$. The hope: an automorphic form built on $6$ might split its
spectrum along the $6N$ wings.

\section*{The data's verdict}
Both attempts return a flat null on the wing, with effect sizes below the noise floor:

\begin{table}[h]
\centering
\begin{tabular}{lll}
\toprule
 & Test A: $a_\ell(E_p)$ & Test B: $a_p$ of level-6 newform \\
\midrule
Wing effect (Cohen's $d$) & $-0.006$ ($a_7$), $-0.004$ ($a_{13}$) & $+0.028$ ($k{=}4$), $+0.043$ ($k{=}6$) \\
Shuffle test & null & $p = 0.625,\ 0.430$ \\
\textbf{Controller of the observable} & $p \bmod \ell$ & none of small modulus \\
\textbf{Variance explained by it} & $\eta^2 = 1.000$ & $\eta^2(p \bmod 12) \approx 0.002$ \\
\bottomrule
\end{tabular}
\end{table}

\noindent
The two controls are mirror images. In Test A the observable is \emph{perfectly} determined by a
congruence: each class $p \bmod \ell$ maps to a single fixed value of $a_\ell$, so $\eta^2 = 1.000$
exactly. In Test B the observable is \emph{immune} to every small congruence: $a_p$ does not track
$p \bmod 12$ ($\eta^2 \approx 0.002$), the signature of Sato--Tate pseudo-randomness. Yet in both, the
6N wing ($p \bmod 3$) explains nothing.

\begin{figure}[h]
\centering
\includegraphics[width=\textwidth]{fig_adversarial.pdf}
\caption{Two opposite mechanisms, one verdict. \emph{Left (Test A):} $a_7(E_p)$ is a step function of
$p \bmod 7$---absolute determinism, $\eta^2 = 1.000$. \emph{Right (Test B):} the level-6 Hecke
eigenvalue $a_p$ shows no structure across $p \bmod 12$ (note the $\pm0.2$ scale)---absolute
pseudo-randomness, $\eta^2 \approx 0.002$. In neither case does the 6N wing reach the observable
($|d| < 0.07$).}
\end{figure}

\section*{Why both fail, and what it confirms}
\textbf{Test A: the fluid is a coordinate.} $a_\ell(E_p)$ depends only on $-p \bmod \ell$, because the
reduction $E_p \bmod \ell$ does. So $a_\ell$ is not a deep fluid sealed inside the CRT lattice---it
\emph{is} a CRT coordinate. Bombarding $p$ at temperature $\ell$ yields exactly $p \bmod \ell$, an axis
orthogonal to the wing $p \bmod 3$ by the Chinese Remainder Theorem. The probe extracts the purest
possible CRT product.

\noindent\textbf{Test B: the gene sits on the wrong primes.} The level $6 = 2\cdot 3$ governs the
\emph{bad} primes $p = 2, 3$; it fixes $a_2, a_3$. For every good prime $p$ (every $6N\pm1$ prime) the
eigenvalue $a_p$ is a Frobenius trace whose value is congruence-immune---and in particular carries no
$p \bmod 3$ information. The ``6 gene'' lives on the level, not on the relation between $a_p$ and the
wing. An object built on $6$ still cannot make its good-prime spectrum feel which wing $p$ lies on.

\noindent\textbf{One verdict.} Orthogonality of parallel coordinates (Test A) and the scale-isolation
of automorphic spectra (Test B) are different physical laws, but they enforce the same boundary: the
control of the 6N wing is locked into elementary number theory and the single $\chi(3)$ Euler channel.
Neither a sharper extraction temperature nor a more sympathetic ambient geometry lets $p \bmod 3$ act
on a good-prime $L$-coefficient. The $\chi(3)$-channel principle survives both attacks; having chosen
the attacks to be as favourable to the wing as we could make them, we regard the boundary as
confirmed.

\paragraph{Archival note.} For Test A: \texttt{Galois\_Fluid\_Isotopes\_Log.csv} (2260 primes),
\texttt{Galois\_fluid\_log.sage}, \texttt{Galois\_fluid\_analysis.py}. For Test B:
\texttt{Langlands\_Plume\_Log.csv} (1334 rows; $k=4,6$), \texttt{Langlands\_plume\_log.sage},
\texttt{Langlands\_plume\_analysis.py}. All archived at
\url{https://github.com/Ruqing1963/6N-chi3-channel}. Each positive/mechanism control certifies its
instrument ($\eta^2 = 1.000$ in A; the explicit $p \bmod 12$ table in B), so both nulls carry the
highest confidence: the instruments are provably working and the wing is provably absent.
\end{document}
